\section{\ENG{End-to-End Runtime Interaction Overview}}
\label{sec:end_to_end_flow}

Η πλήρης λειτουργική ακολουθία ξεκινά από τη διοικητική εισαγωγή ενός μαθήματος. Ο διαχειριστής ανεβάζει το \ENG{ZIP} πακέτο μέσω του \ENG{LMS}, το οποίο χρησιμοποιεί το \ENG{SDK} για να καλέσει το \texttt{\ENG{POST /api/v1/courses/import}}. Ο \ENG{engine} αποθηκεύει το συμπιεσμένο πακέτο στο \ENG{MinIO}, αναλύει το \texttt{\ENG{imsmanifest.xml}} μέσω του \texttt{\ENG{ManifestParser}}, ανιχνεύει το πρότυπο και αποθηκεύει στη \ENG{Postgres} το μάθημα με το αντίστοιχο \ENG{entrypoint} και τα βασικά μεταδεδομένα του. Σε αυτό το στάδιο δεν έχει ακόμη δημιουργηθεί κανένα \ENG{launch} ούτε ενεργή \ENG{runtime} κατάσταση.

Η δεύτερη φάση αφορά την προετοιμασία του εκπαιδευόμενου. Το \ENG{LMS} συγχρονίζει τον χρήστη με το \ENG{engine}, δημιουργεί ή επιβεβαιώνει την εγγραφή του στο μάθημα, και όταν ζητηθεί εκκίνηση καλεί το \texttt{\ENG{POST /api/v1/launches}}. Στο αποθετήριο αναφοράς αυτή η κλήση γίνεται με \texttt{\ENG{forceNewAttempt=true}}, γεγονός που σημαίνει ότι κάθε νέα εκκίνηση δημιουργεί νέα προσπάθεια αντί να αναλαμβάνει αυτόματα συνέχιση ήδη ενεργού \ENG{attempt}. Ο \ENG{engine} δημιουργεί \ENG{attempt}, \ENG{launch} και \ENG{launch token}.

Η τρίτη φάση είναι η μετάβαση στον \ENG{player}. Ο \ENG{browser} φορτώνει το \texttt{\ENG{/launch/:launchId?token=...}}, ο \ENG{player} ανακτά το \ENG{launch context} από το \ENG{engine}, ζητά το \ENG{ZIP} από το \ENG{MinIO}, το αποσυμπιέζει στον τοπικό \ENG{cache} του και αποδίδει \ENG{HTML} με \ENG{iframe} προς το \ENG{entrypoint} του μαθήματος. Στο ίδιο \ENG{HTML} εισάγονται γέφυρες \ENG{SCORM API} στην πλευρά του \ENG{browser}, οι οποίες συλλέγουν \ENG{SetValue}, \ENG{Commit} και \ENG{Terminate} αιτήματα από το \ENG{SCO}. Στο σημείο αυτό ο εκπαιδευόμενος αλληλεπιδρά ήδη με το περιεχόμενο, αλλά η έγκυρη κατάσταση της συνεδρίας εξακολουθεί να βρίσκεται εκτός \ENG{browser}.

Η τέταρτη φάση είναι ο επαναληπτικός κύκλος ενημέρωσης προόδου. Ο \ENG{player} προωθεί τα \ENG{commit} αιτήματα στο \texttt{\ENG{/api/v1/runtime/launches/\{launchId\}/commit}}. Ο \texttt{\ENG{RuntimeOrchestrator}} αναθέτει στον κατάλληλο \ENG{adapter} την ερμηνεία του \ENG{payload}, ενημερώνει το \texttt{\ENG{NormalizedProgress}} και αποθηκεύει το ανανεωμένο \ENG{LaunchRuntimeState} στη \ENG{Redis}. Η \ENG{Postgres} δεν ενημερώνεται σε κάθε \ENG{commit} με πλήρη εγγραφή \ENG{attempt}; ενημερώνεται όταν η κατάσταση συγχρονιστεί προς μόνιμη αποθήκευση κατά τον τερματισμό ή όταν ενεργοποιηθεί ο μηχανισμός για στάσιμες συνεδρίες.

Η πέμπτη και τελική φάση είναι ο τερματισμός και η ανάκτηση του αποτελέσματος. Ο \ENG{player} στέλνει \ENG{terminate} αίτημα στο \ENG{engine}, ο \texttt{\ENG{RuntimeFlushService}} αποθηκεύει στιγμιότυπο, ενημερώνει το \ENG{attempt} με την τελευταία κανονικοποιημένη πρόοδο, κλείνει το \ENG{launch} και καθιστά την κατάσταση πλέον μόνιμα καταγεγραμμένη στο \ENG{Postgres}. Κατόπιν ο \ENG{player} αποστέλλει \texttt{\ENG{postMessage}} προς το \ENG{LMS}, το οποίο ανακατευθύνει τον χρήστη στη σελίδα προόδου και χρησιμοποιεί το \ENG{SDK} για να καλέσει το \texttt{\ENG{GET /api/v1/attempts/\{attemptId\}/progress}}. Η ακολουθία αυτή είναι σημαντική επειδή δείχνει ότι η τελική προβολή προόδου γίνεται ξανά μέσα στο \ENG{LMS}, αλλά πάνω σε έγκυρα δεδομένα αναφοράς που παρήχθησαν από τον \ENG{engine}.

Το Σχήμα~\ref{fig:end_to_end_flow} παρουσιάζει αυτή την πλήρη διαδρομή ως ενιαίο \ENG{sequence}. Η αξία του σχήματος έγκειται στο ότι κάνει ορατές όχι μόνο τις αλληλεπιδράσεις αλλά και τις μεταβολές ευθύνης από φάση σε φάση.

\begin{figure}[H]
  \centering
\begin{chapterfigurebox}
  \begin{tikzpicture}[
      font=\footnotesize,
      node distance=1.8cm and 2.8cm,
      box/.style={draw, rounded corners, fill=gray!5, align=center, text width=3.0cm, minimum height=1.1cm},
      emph/.style={draw, rounded corners, fill=gray!12, align=center, text width=3.2cm, minimum height=1.1cm, font=\footnotesize\bfseries},
      store/.style={draw, cylinder, shape border rotate=90, aspect=0.28, minimum height=1.2cm, minimum width=3.0cm, align=center, fill=gray!5},
      arrow/.style={-Latex, thick}
    ]
    \node[box] (admin) {\ENG{Admin / Learner}};
    \node[box, below=of admin] (lms) {\ENG{LMS}};
    \node[box, below=of lms] (sdk) {\ENG{SDK}};
    \node[emph, below=2.0cm of sdk] (engine) {\ENG{Engine}};
    \node[emph, right=3.4cm of engine] (player) {\ENG{Player}};

    \node[store, below=1.9cm of engine] (redis) {\ENG{Redis}};
    \node[store, below=1.9cm of player] (minio) {\ENG{MinIO}};
    \node[store, below=1.7cm of redis] (postgres) {\ENG{Postgres}};

    \draw[arrow] (admin) -- (lms);
    \draw[arrow] (lms) -- (sdk);
    \draw[arrow] (sdk) -- node[left, align=center]{\ENG{import, launch}\\\ENG{and progress}} (engine);
    \draw[arrow] (engine.west) to[bend left=20] node[left, align=center]{\ENG{launch}\\\ENG{URL}} (lms.east);
    \draw[arrow] (lms.east) to[bend left=18] node[above, align=center]{\ENG{open}\\\ENG{player}} (player.north west);
    \draw[arrow] (player.west) to[bend left=14] node[above, align=center]{\ENG{get}\\\ENG{context}} (engine.east);
    \draw[arrow] (player.west) to[bend right=14] node[below, align=center]{\ENG{commit}\\\ENG{or terminate}} (engine.east);
    \draw[arrow] (player) -- node[right, align=center]{\ENG{download}\\\ENG{ZIP}} (minio);
    \draw[arrow] (engine) -- (redis);
    \draw[arrow] (engine) -- node[left, align=center]{\ENG{persisted}\\\ENG{progress}} (postgres);
  \end{tikzpicture}
\end{chapterfigurebox}
  \caption{Συνολική \ENG{end-to-end} ροή από την εισαγωγή μαθήματος και τη δημιουργία \ENG{launch} έως τη μόνιμη αποθήκευση της προόδου.}
  \label{fig:end_to_end_flow}
\end{figure}

